% =============================================================================
%  ViscoCrowd — Rapport de projet
%  R5.12 Modélisation mathématique · BUT 3 Informatique · 2026-2027
%
%  Compilation :  make report        (depuis la racine du dépôt)
%            ou :  latexmk -pdf main.tex
%
%  Chaque section vit dans son propre fichier sous sections/, pour que trois
%  personnes puissent rédiger en parallèle sans conflit de fusion permanent.
% =============================================================================

\documentclass[11pt,a4paper]{article}

\usepackage[french]{babel}
\usepackage[T1]{fontenc}
\usepackage[utf8]{inputenc}
\usepackage{lmodern}
\usepackage[margin=2.5cm]{geometry}

\usepackage{amsmath,amssymb,amsthm}
\usepackage{graphicx}
\usepackage{booktabs}
\usepackage{siunitx}
\usepackage{caption}
\usepackage{subcaption}
\usepackage{xcolor}
\usepackage[hidelinks]{hyperref}
\usepackage{cleveref}

\graphicspath{{figures/}}

% --- Notations du projet -----------------------------------------------------
% Centraliser les macros évite que trois rédacteurs notent la même grandeur de
% trois façons différentes — un défaut d'homogénéité que le barème sanctionne.
\newcommand{\vect}[1]{\mathbf{#1}}
\newcommand{\pos}{\vect{x}}             % position
\newcommand{\vit}{\vect{v}}             % vitesse
\newcommand{\norm}[1]{\left\lVert #1 \right\rVert}
\newcommand{\proj}[1]{P_{#1}}           % opérateur de projection
\newcommand{\dt}{\Delta t}
\newcommand{\admissible}{\mathcal{K}}   % ensemble des configurations admissibles
\newcommand{\densite}{\rho}
\newcommand{\debit}{q}

\sisetup{locale = FR, per-mode = symbol}

% =============================================================================

\title{%
    \vspace{-1.5cm}
    \textbf{ViscoCrowd}\\[0.3cm]
    \large Extraction aquatique répartie, carrelage mouillé et goulot unique :\\
    effets sur le diagramme fondamental d'une foule en panique
}
\author{%
    Karrouchi Hamza \and Lesueur Shaleena \and Mongrandi Lenny
}
\date{R5.12 — Modélisation mathématique \\ BUT 3 Informatique \\ \today}

\begin{document}

\maketitle

% --- Résumé ------------------------------------------------------------------
% Une demi-page au grand maximum : le problème étudié et l'approche retenue.
\begin{abstract}
% Le resume presente rapidement le probleme etudie et l'approche utilisee.
% Une demi-page au grand maximum.

\textbf{A rediger en fin de projet}, quand les resultats sont connus : un resume
ecrit trop tot promet ce qu'on n'a pas encore mesure.

A couvrir : le phenomene (evacuation de piscine en panique), la question
(deformation du diagramme fondamental par l'enchainement extraction repartie,
zone mouillee, goulot), l'approche (modele particulaire a exclusion dure par
projection), et le resultat principal en une phrase.

\end{abstract}

\tableofcontents
\newpage

% --- Corps du rapport --------------------------------------------------------
% Aucune section ne s'intitule « Développement » : le terme désigne le cœur du
% rapport, qui contient autant de sections que nécessaire.
\section{Introduction}

% L'introduction doit a minima presenter :
%   - le probleme etudie, plus en detail que dans le resume ;
%   - les idees principales de la modelisation ;
%   - le plan suivi dans le rapport.
% Un expose du contexte scientifique sur les approches possibles est un PLUS,
% avec les references bibliographiques correspondantes.

\subsection{Le probleme etudie}

% Evacuation d'urgence d'une piscine couverte. Decrire les trois milieux et
% pourquoi leur enchainement est interessant.

\subsection{Contexte scientifique}

% Les deux grandes familles de modeles de foule : approches continues (la foule
% comme un fluide) et approches particulaires (chaque individu est un objet).
% Au sein des approches particulaires : forces sociales (Helbing) contre
% contraintes dures (Maury et Venel). Dire pourquoi nous retenons la seconde.
% Introduire le diagramme fondamental comme observable de reference.

\subsection{Plan du rapport}

% A rediger en dernier.

\section{Modelisation}

% Le modele mathematique, avant toute discretisation.

\subsection{Etat du systeme}

% Position, vitesse, rayon. Unites.

\subsection{Vitesse desiree et champ de guidage}

% Equation eikonale, et pourquoi le partage rebord/echelles en EMERGE au lieu
% d'etre impose.

\subsection{Les trois milieux}

% Temps de relaxation et acceleration maximale. Justifier physiquement le
% contraste d'adherence entre carrelage mouille et sol sec.

\subsection{Principe d'exclusion}

% L'ensemble admissible, et le fait qu'il N'EST PAS convexe. C'est le point
% delicat du rapport : ne pas le passer sous silence.

\subsection{Franchissements a capacite limitee}

\section{Discretisation numerique}

% Du modele continu au schema numerique.

\subsection{Schema d'Euler semi-implicite}

% Pourquoi semi-implicite plutot qu'explicite : meilleur comportement en temps
% long sur les systemes mecaniques, et c'est lui qui rend possible la
% formulation par minimisation.

\subsection{Formulation par minimisation}

% Un pas de temps vu comme la minimisation d'un compromis entre inertie et
% energie potentielle.

\subsection{Prise en compte des contraintes : une projection}

% Minimiser sous contrainte revient a projeter la position libre sur
% l'ensemble admissible.

\subsection{Linearisation et projection iterative}

% L'ensemble admissible n'etant pas convexe, on linearise autour de la
% configuration courante. Decrire l'algorithme et son critere d'arret.

\subsection{Relecture de la vitesse}

% $\vit^{n+1} = (\pos^{n+1} - \pos^n)/\dt$. C'est la que le choc se produit :
% la vitesse n'est pas transportee, elle est deduite du deplacement reellement
% effectue. Insister, c'est le coeur du schema.

\subsection{Choix du pas de temps}

% Critere : le deplacement par pas doit rester petit devant le rayon.

\section{Implementation}

% Structure du code et choix techniques.

\subsection{Organisation du code}

% Separation moteur / modele / observables, et pourquoi les trajectoires sont
% persistees plutot que les observables calcules au vol.

\subsection{Recherche des voisins}

% Grille de hachage spatial, pour eviter le cout en $N^2$.

\subsection{Reproductibilite}

% Graine unique, manifeste des figures, conteneur Docker.

\section{Validation}

% Une simulation qui tourne n'est pas une simulation juste.
% Cette section doit donner des NOMBRES, pas des impressions visuelles.

\subsection{Validation du moteur sur un cas de reference}

% La bille dans une boite en rotation. Violation de contrainte mesuree a
% \num{2.2e-16} m, independante de $\dt$. Energie strictement non croissante
% en boite fixe.
%
% Figures disponibles : fig01_respect_contrainte, fig02_energie,
% fig03_entrainement, fig04_trajectoire.
%
% Resultat a mettre en avant : le critere d'entrainement solide n'est verifie
% que dans le regime centrifuge ($\omega^2 R / g \gtrsim 1$), et les
% parametres suggeres par le sujet ne s'y trouvent pas. Bon exemple d'analyse
% critique d'un resultat de simulation.

\subsection{Invariants du modele de foule}

% Non-interpenetration, non-sortie du domaine, conservation du nombre d'agents.
% Donner l'ecart mesure, pas seulement une affirmation.

\subsection{Sensibilite a la discretisation}

% $\dt$ divise par 2 puis par 4 ; $N$ croissant ; plusieurs graines.

\subsection{Comparaison a un resultat connu}

% Diagramme fondamental de Weidmann et debit specifique de goulot de la
% litterature.

\subsection{Coherence dimensionnelle}

% Tableau de toutes les grandeurs avec leurs unites et leurs ordres de grandeur.

\section{Interpretation et analyse des resultats}

% Les reponses a la question posee en introduction.

\subsection{Diagramme fondamental de reference}

% Campagne E0 : goulot sec seul.

\subsection{Effet de l'extraction repartie}

% Campagne E1.

\subsection{Effet de l'inertie du carrelage mouille}

% Campagne E2, puis balayage E3. C'est le coeur du resultat : effet tampon ou
% effet inertiel, et de quelle amplitude.

\subsection{Sensibilite aux parametres de scenario}

% Campagne E4.

\section{Conclusion}

% La conclusion doit au moins :
%   - repondre, au moins partiellement, a la problematique posee ;
%   - recapituler les resultats obtenus ;
%   - comporter une analyse critique de ces resultats ;
%   - proposer des pistes d'amelioration.

\subsection{Reponse a la problematique}

\subsection{Analyse critique}

% Limites du modele : 2D, chocs mous uniquement, pas d'heterogeneite
% comportementale fine, calibration des parametres de milieu.

\subsection{Pistes d'amelioration}


% --- Références --------------------------------------------------------------
\bibliographystyle{plain-fr}
\bibliography{references}

% --- Annexes -----------------------------------------------------------------
\appendix
\section{Repartition des taches}

% Tableau indiquant qui a travaille sur quoi et dans quelle proportion.
% La repartition doit etre EQUILIBREE SUR LE FOND : aucun membre ne peut se
% cantonner a la redaction du rapport ou a la preparation de la presentation.
%
% A tenir a jour en continu, en s'appuyant sur l'historique git.

\begin{table}[h]
\centering
\begin{tabular}{lccc}
\toprule
\textbf{Tache} & \textbf{Hamza} & \textbf{Shaleena} & \textbf{Lenny} \\
\midrule
Modelisation mathematique      & & & \\
Schema numerique               & & & \\
Moteur et projection           & & & \\
Champ de guidage               & & & \\
Scenario et geometrie          & & & \\
Observables et campagnes       & & & \\
Validation                     & & & \\
Redaction du rapport           & & & \\
Transparents                   & & & \\
\bottomrule
\end{tabular}
\caption{Repartition des taches au sein du groupe, en pourcentage.}
\end{table}

\section{Outils utilises}

% Les consignes demandent une section indiquant quels outils, notamment d'IA
% generative, ont ete utilises et POUR QUELS USAGES.
%
% Rappel des trois conditions posees par l'enseignant :
%   - l'outil doit permettre de faire plus, ou mieux, que sans lui ;
%   - chaque ligne de code et chaque notion doit pouvoir etre expliquee en
%     detail par les membres du groupe ;
%   - ce que produit une IA doit etre verifie et justifie.
%
% A remplir honnetement et precisement : la soutenance verifie naturellement
% ces conditions.

\subsection{Outils de developpement}

% Python, NumPy, SciPy, Matplotlib, pytest, ruff, Docker, GitHub Actions.

\subsection{Outils d'intelligence artificielle generative}

% Lister : quel outil, pour quelle tache, et comment le resultat a ete verifie.


\end{document}
